\documentclass[11pt,a4paper]{article}
\usepackage[margin=3cm]{geometry}
\usepackage{amsmath,amssymb,amsthm,mathrsfs}
\usepackage[T1]{fontenc}
\usepackage{microtype}
\usepackage[colorlinks=true,linkcolor=blue,citecolor=blue,hypertexnames=false]{hyperref}

\newtheorem{theorem}{Theorem}[section]
\newtheorem{lemma}[theorem]{Lemma}
\newtheorem{proposition}[theorem]{Proposition}
\newtheorem{corollary}[theorem]{Corollary}
\newtheorem{hypothesis}[theorem]{Hypothesis}
\theoremstyle{remark}
\newtheorem{remark}[theorem]{Remark}

\makeatletter
\renewcommand{\@seccntformat}[1]{\makebox[2.1em][l]{\csname the#1\endcsname}}
\makeatother

\newcommand{\R}{\mathbb R}
\newcommand{\A}{\mathfrak A}
\newcommand{\G}{\mathcal G}
\newcommand{\Hh}{\mathcal H_0}
\newcommand{\Om}{\Omega}
\DeclareMathOperator{\supp}{supp}
\DeclareMathOperator{\dist}{dist}
\DeclareMathOperator{\spec}{spec}
\DeclareMathOperator{\gap}{gap}

\title{M4a: a conditional continuum LPPL theorem\\
for removal of the spatial cutoff in $P(\varphi)_2$}
\author{}
\date{}

\begin{document}
\maketitle

\begin{abstract}
This note proves the continuum Cauchy step of the spatial-cutoff programme in an explicitly stated
operator-localization framework and under two analytic hypotheses: a uniform Hamiltonian gap and
a uniform local second moment of the interaction density.  Exact finite propagation for the
\emph{actual cutoff dynamics} and the required ground-state operator domains are proved from the
standard cutoff construction.  Affine-path differentiability and endpoint continuity are proved
from the Feynman--Kac--Nelson formula and the assumed gap.  No bounded quasi-adiabatic generator is
used.
Instead, differentiating the state gives a connected correlation with one unbounded argument.
An exact spectral filter and a gap-plus-locality estimate reduce this correlation to a summable
one-dimensional cell expansion.

The conclusion is norm convergence of the full cutoff net on every local algebra, with a rate
faster than every inverse power of the distance to the cutoff, local normality, translation
invariance, $\mathbb Z_2$-invariance for even $P$, and the algebraic ground-state property proved
in M0.  The theorem does \emph{not} prove the two remaining hypotheses, uniqueness of all ground states,
the sharp rate, or uniformity in a wavelet resolution.  Those points are separated from the
conditional theorem rather than hidden inside its proof.
\end{abstract}

\section{Framework and hypotheses}\label{sec:framework}

For every bounded open $U\subset\R$, let $\A(U)$ be the von Neumann algebra generated by the
time-zero Weyl operators whose test functions have support in $U$.  Let $O\subset\R$ be a bounded
open interval and $O_r:=\{x\in\R:\dist(x,O)<r\}$.  We use the time-zero local net
$U\mapsto\A(U)\subset B(\Hh)$ and its norm-closed quasi-local algebra
$\A=\overline{\bigcup_O\A(O)}^{\|\cdot\|}$.  The net is isotonic and local:
\[
 O_1\cap O_2=\varnothing\quad\Longrightarrow\quad
 [\A(O_1),\A(O_2)]=0.
\]
For real $h\in C_c^\infty(\R)$, write
\[
 V(h):=\int_{\R}h(x):P(\varphi_0(x)):\,dx .
\]
The operator-localization framework used below is the following.  For every real
$h\in C_c^\infty(\R)$, $V(h)$ is a specified self-adjoint operator; it is affiliated with $\A(U)$
for every bounded open $U$ containing $\supp h$.  If $h=\sum_{j=1}^N h_j$ and
$\xi\in D(V(h))\cap\bigcap_{j=1}^ND(V(h_j))$, then
\[
 V(h)\xi=\sum_{j=1}^N V(h_j)\xi.
\]
These assertions are part of the standard operator-localization framework.

Let
\[
 \G:=\{g\in C_c^\infty(\R):0\le g\le1,\ g=1
 \text{ on a nonempty bounded interval}\},
\]
ordered by $g\preccurlyeq g'$ if $\{g=1\}^\circ\subseteq\{g'=1\}^\circ$.  Put
\[
 K(g):=H_0+V(g),\qquad E(g):=\inf\spec K(g),\qquad H(g):=K(g)-E(g).
\]
For each $g\in\G$, assume that $0$ is a simple eigenvalue of $H(g)$ with normalized eigenvector
$\Om_g$.  Set
\[
 \omega_g(A):=\langle\Om_g,A\Om_g\rangle,
 \qquad \beta_t^g(A):=e^{itH(g)}Ae^{-itH(g)}.
\]
Inner products are linear in the second argument.

This order is directed.  Indeed, for $g_1,g_2\in\G$, choose a bounded interval $I$ containing
$\{g_1=1\}\cup\{g_2=1\}$ and a function $g_3\in C_c^\infty(\R;[0,1])$ equal to $1$ on a
neighbourhood of $\overline I$; then $g_1,g_2\preccurlyeq g_3$.  Similarly, for every bounded
interval $O$ and $d<\infty$, one may choose $g_0\in\G$ equal to $1$ on $O_d$, and every
$g\succcurlyeq g_0$ is also equal to $1$ on $O_d$.

The following are the exact hypotheses used below.

\begin{hypothesis}[uniform gap, G]\label{hyp:G}
There is $\gamma>0$ such that, for every $g\in\G$,
\[
 \spec H(g)\subseteq\{0\}\cup[\gamma,\infty).
\]
\end{hypothesis}

\begin{hypothesis}[uniform local interaction moment, M]\label{hyp:M}
There is $M<\infty$ such that, for every $g\in\G$ and every real
$h\in C_c^\infty(\R)$ with $\|h\|_\infty\le1$ whose support is contained in an interval of length
at most $2$,
\[
 \|V(h)\Om_g\|\le M.
\]
\end{hypothesis}

\begin{proposition}[domain smoothing and cutoff propagation]\label{prop:standard}
The standard spatial-cutoff $P(\varphi)_2$ construction has the following two consequences.
\begin{enumerate}
\item For every $g\in\G$ and every real $h\in C_c^\infty(\R)$,
$\Om_g\in D(V(h))$.
\item For every $g\in\G$, every bounded interval $O$, and every $t\in\R$,
\[
 \beta_t^g(\A(O))\subseteq\A(O_{|t|}).
\]
The propagation-speed bound is one and is independent of $g$.
\end{enumerate}
\end{proposition}

\begin{proof}
For the first assertion, the Glimm--Jaffe semigroup-smoothing theorem
\cite[Ch.~7, Thm.~7.3 and Lem.~7.4]{GJExp} gives, for $t>0$ and every
compactly smeared Wick polynomial $V(h)$,
\[
 \operatorname{Ran}e^{-tK(g)}\subseteq D(V(h)).
\]
Since $K(g)\Om_g=E(g)\Om_g$,
$\Om_g=e^{tE(g)}e^{-tK(g)}\Om_g\in D(V(h))$.

For propagation, the interaction automorphism
$\beta_t^{I,g}=\operatorname{Ad}e^{itV(g)}$ has speed zero.  Indeed, for an observable generated
by Weyl operators supported in a compact $B\Subset O$, choose $\chi\in C_c^\infty(O)$ equal to
one near $B$.  The strongly commuting decomposition
\[
 V(g)=V(\chi g)+V((1-\chi)g)
\]
has $e^{itV(\chi g)}\in\A(O)$ and
$e^{itV((1-\chi)g)}\in\A(B)'$, whence
$\beta_t^{I,g}(A)\in\A(O)$.  Strong density extends this to every $A\in\A(O)$.
The free dynamics $\beta_t^0=\operatorname{Ad}e^{itH_0}$ has speed one by the finite-propagation
property of the Klein--Gordon equation.  Therefore the $n$th Trotter approximant
\[
 \bigl(\beta_{t/n}^0\circ\beta_{t/n}^{I,g}\bigr)^n(A)
\]
belongs to $\A(O_{|t|})$.  The Trotter product formula gives strong convergence of these
approximants to $\beta_t^g(A)$, and the von Neumann algebra $\A(O_{|t|})$ is strongly closed.
This is the standard ``free speed one + interaction speed zero'' proof
\cite[Ch.~4, Prop.~4.1.1 and Thm.~4.1.2]{GJExp}.
\end{proof}

\begin{proposition}[affine-path regularity from FKN]\label{prop:FKNpath}
Let $g,g'\in\G$ be equal to one on a common nonempty interval, set
$g_s=(1-s)g+sg'$, $v=g'-g$, and abbreviate
\[
 K_s=K(g_s),\quad E_s=E(g_s),\quad H_s=K_s-E_s,\quad
 P_s=|\Om_s\rangle\langle\Om_s|,\quad Q_s=1-P_s.
\]
Under Hypothesis~\ref{hyp:G}, $s\mapsto P_s$ is norm-continuous on $[0,1]$ and $C^1$ on
$(0,1)$.  Locally on $(0,1)$, normalized ground-state vectors can be chosen $C^1$ in Hilbert norm
and in the parallel gauge.  They then obey
\begin{equation}\label{eq:FKNderivatives}
 \dot E_s=\langle\Om_s,V(v)\Om_s\rangle,
 \qquad
 \dot\Om_s=-H_s^{-1}Q_sV(v)\Om_s.                         \tag{FKN}
\end{equation}
In particular, $s\mapsto\omega_{g_s}(A)=\operatorname{Tr}(P_sA)$ is continuous on $[0,1]$ for
every $A\in B(\Hh)$.
\end{proposition}

\begin{proof}
Let $m_0>0$ be the free mass and fix $T>0$.  On free Euclidean path space, with time-slice maps
$J_0,J_T$, define
\[
 \mathcal U_s=\int_0^T\!\int_{\R}g_s(x):P(\varphi(t,x)):\,dx\,dt,
 \qquad w_s=e^{-\mathcal U_s},
 \qquad \mathcal U_v=\mathcal U_1-\mathcal U_0.
\]
For bounded time-zero functions $f,h$, the Feynman--Kac--Nelson formula is
\begin{equation}\label{eq:FKN}
 \langle f,e^{-TK_s}h\rangle
 =\int \overline{J_Tf}\,(J_0h)w_s\,d\mu_0.                 \tag{1.1}
\end{equation}
Put
\[
 q_T=\frac{2}{1-e^{-m_0T}},\qquad r_T=\frac{q_T}{q_T-1}.
\]
Here is the full derivation of the time-slice product estimate.  Let $\nu_0$ be the time-zero
Gaussian marginal of $\mu_0$, and let $\mathsf P_T=e^{-TH_0}$ be its Markov semigroup.  Since
\[
 r_T=\frac{2}{1+e^{-m_0T}},\qquad
 p_T:=\frac2{r_T}=1+e^{-m_0T},\qquad p_T'=1+e^{m_0T},
\]
$p_T'$ is conjugate to $p_T$ and
$(p_T'-1)/(p_T-1)=e^{2m_0T}$.  Nelson hypercontractivity gives
$\|\mathsf P_TF\|_{p_T'}\le\|F\|_{p_T}$.  Hence the Markov property and H\"older's inequality
give
\begin{align*}
 \|(J_Tf)(J_0h)\|_{r_T}^{r_T}
 &=\int |h|^{r_T}\,\mathsf P_T(|f|^{r_T})\,d\nu_0\\
 &\le\||h|^{r_T}\|_{p_T}
       \|\mathsf P_T(|f|^{r_T})\|_{p_T'}\\
 &\le\||h|^{r_T}\|_{p_T}\||f|^{r_T}\|_{p_T}
 =\|h\|_2^{r_T}\|f\|_2^{r_T}.
\end{align*}
Taking the $r_T$th root yields
\begin{equation}\label{eq:timeslice}
 \|(J_Tf)(J_0h)\|_{r_T}\le\|f\|_2\|h\|_2.                \tag{1.2}
\end{equation}
The underlying hypercontractivity and transfer-matrix facts are
\cite[Thms.~V.10 and V.12]{Simon}.

The spacetime-smeared Wick polynomial $\mathcal U_v$ belongs to $L^p(\mu_0)$ for every finite
$p$.  Since $g,g'\ge0$ and $P$ is semibounded, the exponential estimate
\cite[Thm.~V.7]{Simon} gives
\[
 e^{-a\mathcal U_0},e^{-a\mathcal U_1}\in L^1(\mu_0)
 \qquad(a<\infty).
\]
Because $\mathcal U_s=(1-s)\mathcal U_0+s\mathcal U_1$, H\"older's inequality gives
\begin{equation}\label{eq:FKNmoments}
 \|w_s\|_a^a
 \le
 \left(\int e^{-a\mathcal U_0}\,d\mu_0\right)^{1-s}
 \left(\int e^{-a\mathcal U_1}\,d\mu_0\right)^s.          \tag{1.3}
\end{equation}
Thus $\sup_{0\le s\le1}\|w_s\|_a<\infty$ for each finite $a$.

If $s_n\to s$, then $\mathcal U_{s_n}\to\mathcal U_s$ in every finite $L^p$, so
$w_{s_n}\to w_s$ in probability.  Applying \eqref{eq:FKNmoments} with an exponent strictly
larger than $a$ makes $(|w_{s_n}|^a)_n$ uniformly integrable.  Vitali's theorem therefore yields
\begin{equation}\label{eq:FKNweightcontinuity}
 \|w_{s_n}-w_s\|_a\longrightarrow0                         \tag{1.4}
\end{equation}
for every finite $a$.  Fix $q=q_T$ and choose finite $p,a$ with
$q^{-1}=p^{-1}+a^{-1}$.  Whenever $s,s+h\in[0,1]$, the scalar fundamental theorem of calculus
gives, pointwise in path space,
\[
 \frac{w_{s+h}-w_s}{h}
 =-\mathcal U_v\int_0^1w_{s+\theta h}\,d\theta.
\]
Consequently, for $0<s<1$,
\[
 \left\|\frac{w_{s+h}-w_s}{h}+\mathcal U_vw_s\right\|_q
 \le \|\mathcal U_v\|_p
       \sup_{0\le\theta\le1}\|w_{s+\theta h}-w_s\|_a
 \longrightarrow0.
\]
Thus $s\mapsto w_s$ is continuous from $[0,1]$ to $L^{q_T}$ and is $C^1$ from $(0,1)$ to
$L^{q_T}$, with derivative $-\mathcal U_vw_s$.  The derivative is continuous because
\[
 \|\mathcal U_v(w_s-w_r)\|_{q_T}
 \le\|\mathcal U_v\|_p\|w_s-w_r\|_a\longrightarrow0.
\]
Equations
\eqref{eq:FKN}--\eqref{eq:timeslice}, first for bounded $f,h$ and then by $L^2$ density, imply
\[
 \|e^{-TK_s}-e^{-TK_r}\|\le\|w_s-w_r\|_{q_T},
\]
and the identical estimate for difference quotients proves that
$s\mapsto L_s:=e^{-TK_s}$ is norm-continuous on $[0,1]$ and norm-$C^1$ on $(0,1)$.

The largest spectral value of $L_s$ is the simple eigenvalue
$\lambda_s=e^{-TE_s}$.  Hypothesis~\ref{hyp:G} gives
\[
 \spec L_s\setminus\{\lambda_s\}
 \subset[0,e^{-T\gamma}\lambda_s].
\]
Since $\lambda_s=\|L_s\|$, it is continuous; hence the displayed relative separation is locally
bounded below by a positive absolute distance.  The Riesz integral around $\lambda_s$ now proves
that $P_s$ is norm-continuous on $[0,1]$ and $C^1$ on $(0,1)$.  Near any
$s_0\in(0,1)$,
\[
 \widetilde\Om_s=\frac{P_s\Om_{s_0}}{\|P_s\Om_{s_0}\|}
\]
is $C^1$ after shrinking the neighbourhood.  A $C^1$ phase change puts it in the parallel gauge
$\langle\Om_s,\dot\Om_s\rangle=0$.  The identity
$\lambda_s=\langle\Om_s,L_s\Om_s\rangle$ then shows explicitly that $\lambda_s$, and hence
$E_s=-T^{-1}\log\lambda_s$, is $C^1$ on $(0,1)$.

It remains to identify the derivative.  The cutoff Hamiltonians have the common operator core
$\mathscr C=C^\infty(H_0)$ and, on that core,
\[
 K_s\xi=K_{s_0}\xi+(s-s_0)V(v)\xi,
 \qquad \xi\in\mathscr C,
\]
by \cite[Thm.~V.12(a)]{Simon}.  For fixed $\xi\in\mathscr C$, differentiate
$\langle K_s\xi,\Om_s\rangle=E_s\langle\xi,\Om_s\rangle$.  Here
$E_s=-T^{-1}\log\lambda_s$ is $C^1$, and Proposition~\ref{prop:standard} gives
$\Om_s\in D(V(v))$.  Self-adjointness of $V(v)$ therefore yields
\[
 \langle H_s\xi,\dot\Om_s\rangle
 =\langle\xi,(\dot E_s-V(v))\Om_s\rangle
 \qquad(\xi\in\mathscr C).
\]
Since $\mathscr C$ is a core for $H_s$, the definition of the adjoint implies
$\dot\Om_s\in D(H_s)$ and
\[
 H_s\dot\Om_s=(\dot E_s-V(v))\Om_s.
\]
Taking the inner product with $\Om_s$ gives the first identity in
\eqref{eq:FKNderivatives}; the parallel gauge and
$\|H_s^{-1}Q_s\|\le\gamma^{-1}$ give the second.  Finally,
\[
 |\operatorname{Tr}((P_s-P_r)A)|
 \le\|P_s-P_r\|_1\|A\|
 \le2\|P_s-P_r\|\,\|A\|,
\]
because $P_s-P_r$ has rank at most two.  This proves the endpoint assertion.
\end{proof}

For the final ground-state identification and invariances, we also use the already established
facts
\begin{align}
 &\alpha_t(\A(O))\subseteq\A(O_{|t|}),\qquad
 \alpha_t(A)=\beta_t^g(A)\quad\text{if }g=1\text{ on }O_{|t|},
 \label{eq:physical}\tag{P}\displaybreak[0]\\
 &\omega_{\tau_ag}=\omega_g\circ\tau_{-a},
 \label{eq:translation}\tag{T}
\end{align}
where $(\tau_ag)(x)=g(x-a)$ and the same symbol denotes the spatial automorphism of $\A$.

\section{An exact spectral filter}\label{sec:filter}

The proof needs an $L^1$ real-time representation of $H^{-1}$ on the gapped subspace.  The
following construction gives more decay than is needed for convergence.

\begin{lemma}[exact filter]\label{lem:filter}
For every $\gamma>0$ there is $W_\gamma\in L^1(\R)$ such that, with
\[
 \widehat W_\gamma(E):=\int_\R W_\gamma(t)e^{itE}\,dt,
\]
one has
\[
 \widehat W_\gamma(0)=0,
 \qquad \widehat W_\gamma(E)=\frac1E\quad\text{for }|E|\ge\gamma.
\]
Moreover, for every integer $q\ge0$,
\[
 \int_\R(1+|t|)^q|W_\gamma(t)|\,dt<\infty.
\]
Consequently, if $H\ge0$ and $\spec H\subseteq\{0\}\cup[\gamma,\infty)$ with
$Q=1-P_{\ker H}$, then
\[
 H^{-1}Q=\int_\R W_\gamma(t)e^{itH}Q\,dt
\]
as a norm-convergent Bochner integral.
\end{lemma}

\begin{proof}
Choose an even $\chi\in C^\infty(\R;[0,1])$ with
$\chi(E)=0$ for $|E|\le\gamma/2$ and $\chi(E)=1$ for $|E|\ge\gamma$.  Define
\[
 \widehat W_\gamma(E)=
 \begin{cases}\chi(E)/E,&E\ne0,\\0,&E=0.\end{cases}
\]
This is a smooth odd function.  For every $k\ge1$, its $k$th derivative is integrable: it is
smooth on a compact transition region and is a constant multiple of $E^{-k-1}$ outside that
region.

Define the inverse Fourier transform first as an oscillatory integral,
\[
 W_\gamma(t):=\frac1{2\pi}\int_\R e^{-itE}\widehat W_\gamma(E)\,dE.
\]
For $t\ne0$, oddness gives
\[
 W_\gamma(t)=-\frac{i}{\pi}\int_0^\infty
 \sin(tE)\frac{\chi(E)}E\,dE.
\]
This is uniformly bounded for $0<|t|\le1$: the integral over $[0,\gamma]$ is bounded absolutely,
and after the substitution $u=|t|E$ the integral over $[\gamma,\infty)$ is a tail of
$\int_0^\infty\sin u/u\,du$, whose partial integrals are uniformly bounded by Dirichlet's test.
For $|t|\ge1$, integration by parts $k$ times in the oscillatory integral gives
\[
 |W_\gamma(t)|\le\frac{\|\widehat W_\gamma^{(k)}\|_{L^1}}{2\pi|t|^k}
 \qquad(k\ge1),
\]
because every boundary term vanishes.  Taking $k>q+1$ proves the moment assertion and in
particular $W_\gamma\in L^1$.  Fourier inversion in tempered distributions, followed by
continuity of the Fourier transform of an $L^1$ function, gives
$\mathcal F W_\gamma=\widehat W_\gamma$ pointwise.

The spectral theorem now yields
\[
 \int_\R W_\gamma(t)e^{itH}Q\,dt
 =\widehat W_\gamma(H)Q=H^{-1}Q.
\]
The integral converges in operator norm because $W_\gamma\in L^1$ and
$\|e^{itH}Q\|\le1$.
\end{proof}

Put
\[
 T_\gamma(r):=\int_{|t|>r}|W_\gamma(t)|\,dt.
\]
Lemma~\ref{lem:filter} implies, for every $q\ge0$,
\begin{equation}\label{eq:filtertail}
 T_\gamma(r)\le C_{q,\gamma}(1+r)^{-q}.
\end{equation}

\section{Clustering with one unbounded local observable}\label{sec:cluster}

We first isolate the scalar spectral-separation estimate used in M3.

\begin{lemma}[spectral separation]\label{lem:spectral}
Let
\[
 F(t)=\int_{[\gamma,\infty)}e^{itE}\,d\rho(E),\qquad
 G(t)=\int_{(-\infty,-\gamma]}e^{itE}\,d\sigma(E),
\]
where $\rho$ and $\sigma$ are complex measures with
$|\rho|(\R),|\sigma|(\R)\le M_0$.  If $F(t)=G(t)$ for $|t|<R$ and $\gamma R\ge1$, then
\[
 |F(0)|\le2M_0e^{-\gamma R/2}.
\]
\end{lemma}

\begin{proof}
For $a,\varepsilon>0$, set
\[
 f_{a,\varepsilon}(t):=\frac1{2\pi i}\frac{e^{-t^2/(2a)}}{t-i\varepsilon}.
\]
Residue calculus followed by convolution with the Fourier transform of the Gaussian gives
\[
 \check f_{a,\varepsilon}(E):=\int_\R f_{a,\varepsilon}(t)e^{itE}\,dt
 =\int_0^\infty\sqrt{\frac a{2\pi}}
 e^{-a(E-E')^2/2}e^{-\varepsilon E'}\,dE'.
\]
Hence $0\le\check f_{a,\varepsilon}\le1$ and
$\check f_{a,\varepsilon}(E)\to\Phi(\sqrt aE)$ as $\varepsilon\downarrow0$, where $\Phi$ is the
standard normal distribution function.  Also
\[
 \int_{|t|\ge R}|f_{a,\varepsilon}(t)|\,dt
 \le\frac a{\pi R^2}e^{-R^2/(2a)}.
\]
Indeed, $|t-i\varepsilon|\ge|t|$ and
$\int_R^\infty e^{-t^2/(2a)}dt\le(a/R)e^{-R^2/(2a)}$.

Using Fubini, the equality $F=G$ on $(-R,R)$, and the total-variation bounds, then letting
$\varepsilon\downarrow0$, gives
\[
 \left|\int\Phi(\sqrt aE)\,d\rho(E)
       -\int\Phi(\sqrt aE)\,d\sigma(E)\right|
 \le \frac{2M_0a}{\pi R^2}e^{-R^2/(2a)}.
\]
For $x\ge0$, $1-\Phi(x)\le\tfrac12e^{-x^2/2}$ and
$\Phi(-x)\le\tfrac12e^{-x^2/2}$.  The support conditions therefore imply
\[
 |F(0)|\le
 M_0e^{-a\gamma^2/2}
 +\frac{2M_0a}{\pi R^2}e^{-R^2/(2a)}.
\]
Taking $a=R/\gamma$ and using $\gamma R\ge1$ gives
\[
 |F(0)|\le M_0e^{-\gamma R/2}
 \left(1+\frac{2}{\pi\gamma R}\right)
 \le2M_0e^{-\gamma R/2}.
\]
\end{proof}

\begin{lemma}[one unbounded argument]\label{lem:unboundedcluster}
Let $H\ge0$ have normalized simple ground state $\Om$ and
$\spec H\subseteq\{0\}\cup[\gamma,\infty)$.  Suppose its dynamics
$\beta_t=\operatorname{Ad}e^{itH}$ satisfies exact finite propagation.
Let $B\subset\R$ be compact, let $X=X^*$ be affiliated with $\A(U)$ for every bounded open
neighbourhood $U$ of $B$, let $\Om\in D(X)$, and let $A\in\A(O)$, where $O$ is a bounded open
interval.  If $R:=\dist(B,O)>0$ and $\gamma R\ge1$, then
\begin{equation}\label{eq:unboundedcluster}
 \left|\langle X\Om,A\Om\rangle
 -\langle X\Om,\Om\rangle\langle\Om,A\Om\rangle\right|
 \le2\|X\Om\|\,\|A\|e^{-\gamma R/2}.
\end{equation}
\end{lemma}

\begin{proof}
Let $P_0=|\Om\rangle\langle\Om|$ and define
\[
 F(t):=\langle X\Om,(e^{itH}-P_0)A\Om\rangle,
 \qquad
 G(t):=\langle A^*\Om,(e^{-itH}-P_0)X\Om\rangle.
\]
The spectral theorem represents $F$ and $G$ by complex measures supported in
$[\gamma,\infty)$ and $(-\infty,-\gamma]$, respectively, with total variations at most
$M_0:=\|X\Om\|\|A\|$.

Fix $|t|<R$ and choose a bounded open neighbourhood $U$ of $B$ disjoint from $O_{|t|}$.
Finite propagation gives $\beta_t(A)\in\A(O_{|t|})\subseteq\A(U)'$.  Affiliation of $X$ with
$\A(U)$ implies that every element of $\A(U)'$ preserves $D(X)$ and commutes with $X$ there.
Since $e^{-itH}\Om=\Om$,
\[
 \langle X\Om,e^{itH}A\Om\rangle
 =\langle\Om,X\beta_t(A)\Om\rangle
 =\langle\Om,\beta_t(A)X\Om\rangle
 =\langle A^*\Om,e^{-itH}X\Om\rangle.
\]
The two $P_0$ terms are also equal: both are
$\langle X\Om,\Om\rangle\langle\Om,A\Om\rangle$, because
$\langle X\Om,\Om\rangle\in\R$.  Hence $F(t)=G(t)$ for $|t|<R$.
Lemma~\ref{lem:spectral} applies, and $F(0)$ is the expression on the left of
\eqref{eq:unboundedcluster}.
\end{proof}

\section{Differentiating the ground-state expectation}\label{sec:derivative}

\begin{lemma}[state derivative]\label{lem:derivative}
Assume Hypothesis~\ref{hyp:G}.  Let $g,g'\in\G$ be equal to one on a common nonempty interval,
put $g_s=(1-s)g+sg'$, and, for $0<s<1$, write
\[
 V:=V(g'-g),\quad E_s:=E(g_s),\quad H_s:=H(g_s),\quad
 P_s:=|\Om_s\rangle\langle\Om_s|,\quad Q_s:=1-P_s.
\]
Then
\begin{align}
 \dot E_s&=\langle\Om_s,V\Om_s\rangle,\label{eq:HF}\\
 \dot\Om_s&=-H_s^{-1}Q_sV\Om_s,\label{eq:eigenderivative}
\end{align}
and, for every self-adjoint $A\in B(\Hh)$,
\begin{equation}\label{eq:statederivative}
 \frac d{ds}\omega_{g_s}(A)
 =-2\operatorname{Re}\int_\R W_\gamma(t)
 \Bigl[
 \langle V\Om_s,\beta_t^{g_s}(A)\Om_s\rangle
 -\langle V\Om_s,\Om_s\rangle\omega_{g_s}(A)
 \Bigr]dt.
\end{equation}
The integral is absolutely convergent.
\end{lemma}

\begin{proof}
The first two identities are exactly Proposition~\ref{prop:FKNpath}; in particular all products
with the unbounded operator $V$ are defined by Proposition~\ref{prop:standard}.

For self-adjoint $A$,
\[
 \frac d{ds}\omega_{g_s}(A)
 =2\operatorname{Re}\langle\dot\Om_s,A\Om_s\rangle
 =-2\operatorname{Re}\langle V\Om_s,H_s^{-1}Q_sA\Om_s\rangle.
\]
Insert Lemma~\ref{lem:filter}.  Since
$e^{itH_s}A\Om_s=\beta_t^{g_s}(A)\Om_s$ and
$P_sA\Om_s=\omega_{g_s}(A)\Om_s$, this is exactly
\eqref{eq:statederivative}.  Absolute convergence follows from
\[
 \left|\langle V\Om_s,e^{itH_s}Q_sA\Om_s\rangle\right|
 \le\|V\Om_s\|\|A\|
\]
and $W_\gamma\in L^1$.
\end{proof}

\section{The cell estimate and the Cauchy theorem}\label{sec:cauchy}

Choose $\rho\in C_c^\infty((-1,1))$ with $0\le\rho\le1$ and
$\rho>0$ on $[-\tfrac12,\tfrac12]$, put
\[
 D(x):=\sum_{k\in\mathbb Z}\rho(x-k),\qquad
 \eta_j(x):=\frac{\rho(x-j)}{D(x)}.
\]
The sum defining $D$ is locally finite, $D(x)>0$, and $D$ is smooth and one-periodic.  Hence
$(\eta_j)_{j\in\mathbb Z}$ is a smooth partition of unity satisfying
\[
 0\le\eta_j\le1,\qquad \supp\eta_j\subset(j-1,j+1),\qquad
 \sum_{j\in\mathbb Z}\eta_j=1.
\]
For $v=g'-g$, set $h_j=\eta_jv$ and $V_j=V(h_j)$.  Only finitely many $h_j$ are nonzero,
$V(v)=\sum_jV_j$, and Hypothesis~\ref{hyp:M} gives
\begin{equation}\label{eq:cellmoment}
 \|V_j\Om_s\|\le M
 \qquad(0\le s\le1).
\end{equation}

\begin{proposition}[one path]\label{prop:path}
Assume Hypotheses~\ref{hyp:G} and \ref{hyp:M}.  Let $O$ be a bounded interval and suppose
$g=g'=1$ on $O_d$, where $d\ge d_0:=4+2/\gamma$.  Then
\begin{equation}\label{eq:pathbound}
 \left\|(\omega_g-\omega_{g'})|_{\A(O)}\right\|
 \le 32 M\sum_{n\ge\lfloor d\rfloor}
 \left(\|W_\gamma\|_1e^{-\gamma n/4}+T_\gamma(n/2)\right),
\end{equation}
\end{proposition}

\begin{proof}
It is enough to take self-adjoint $A\in\A(O)$ with $\|A\|\le1$, because
$\omega_g-\omega_{g'}$ is hermitian.  Since $v=0$ on $O_d$, every nonzero $h_j$ is supported at
distance at least $d$ from $O$.  Put $B_j=\supp h_j$ and $r_j=\dist(B_j,O)$.

For $0<s<1$, the inequalities $0\le g,g'\le1$ imply
\[
 g_s(x)=1\quad\Longleftrightarrow\quad g(x)=g'(x)=1.
\]
Thus $g_s\in\G$ and $g_s=1$ on $O_d$ for all $s\in[0,1]$, so both hypotheses apply along the
entire path.

For fixed $s,j$, define the vector correlation
\[
 C_{j,s}(t):=
 \langle V_j\Om_s,\beta_t^{g_s}(A)\Om_s\rangle
 -\langle V_j\Om_s,\Om_s\rangle\omega_{g_s}(A).
\]
If $|t|\le r_j/2$, Proposition~\ref{prop:standard} gives
$\beta_t^{g_s}(A)\in\A(O_{|t|})$, whose distance from $B_j$ is at least $r_j/2$.  Moreover,
$\gamma r_j/2\ge1$ because $r_j\ge d\ge2/\gamma$.  Lemma~\ref{lem:unboundedcluster}, applied to
$X=V_j$ and
$\beta_t^{g_s}(A)$, together with \eqref{eq:cellmoment}, yields
\begin{equation}\label{eq:nearbound}
 |C_{j,s}(t)|\le2M e^{-\gamma r_j/4}.
\end{equation}
For arbitrary $t$, Cauchy--Schwarz gives the uniform bound
\begin{equation}\label{eq:farbound}
 |C_{j,s}(t)|
 \le |\langle V_j\Om_s,\beta_t^{g_s}(A)\Om_s\rangle|
     +|\langle V_j\Om_s,\Om_s\rangle|\,|\omega_{g_s}(A)|
 \le2M.
\end{equation}
Substituting $V=\sum_jV_j$ into \eqref{eq:statederivative}, splitting the $t$-integral at
$|t|=r_j/2$, and using \eqref{eq:nearbound}--\eqref{eq:farbound}, gives
\[
 \left|\frac d{ds}\omega_{g_s}(A)\right|
 \le4M\sum_j
 \left(\|W_\gamma\|_1e^{-\gamma r_j/4}+T_\gamma(r_j/2)\right).
\]

Write $O=(a,b)$.  Because $d\ge4$ and $B_j\subset[j-1,j+1]$, a nonempty $B_j$ cannot have points
on both sides of $O_d$.  If $B_j$ is on the right and $r_j\in[n,n+1)$, compactness gives
$x_j\in B_j$ with $x_j-b=r_j$.  Since $|j-x_j|\le1$, necessarily
$j\in[b+n-1,b+n+2]$, an interval of length $3$ containing at most four integers.  The same count
holds on the left.  Consequently at most eight indices occur in each shell
$r_j\in[n,n+1)$.  Both functions
$r\mapsto e^{-\gamma r/4}$ and $r\mapsto T_\gamma(r/2)$ are decreasing.  Since $r_j\ge d$, the
last display is therefore at most
\[
 32M\sum_{n\ge\lfloor d\rfloor}
 \left(\|W_\gamma\|_1e^{-\gamma n/4}+T_\gamma(n/2)\right).
\]
For $0<\varepsilon<1/2$, integrate the last derivative bound over
$[\varepsilon,1-\varepsilon]$.  Proposition~\ref{prop:FKNpath} gives norm continuity of $P_s$ at
$s=0,1$, hence continuity of $\omega_{g_s}(A)=\operatorname{Tr}(P_sA)$ there.  Letting
$\varepsilon\downarrow0$ and taking the supremum over the self-adjoint unit ball proves the
proposition.
\end{proof}

Define
\begin{equation}\label{eq:Psi}
 \Psi_\gamma(d):=32M\sum_{n\ge\lfloor d\rfloor}
 \left(\|W_\gamma\|_1e^{-\gamma n/4}+T_\gamma(n/2)\right)
 \quad(d\ge d_0),
\end{equation}
and set $\Psi_\gamma(d):=\max\{2,\Psi_\gamma(d_0)\}$ for $0\le d<d_0$.  This is a bounded
decreasing function; the value $2$ gives the trivial bound for the norm-distance of two states
when Proposition~\ref{prop:path} is unavailable.

\begin{lemma}[decay of the comparison function]\label{lem:Psitail}
$\Psi_\gamma(d)\to0$ as $d\to\infty$.  More precisely, for every integer $N\ge1$ there is
$C_{N,\gamma,M}<\infty$ such that
\[
 \Psi_\gamma(d)\le C_{N,\gamma,M}(1+d)^{-N}.
\]
\end{lemma}

\begin{proof}
The exponential series in \eqref{eq:Psi} has an exponentially decaying tail.  For the second
series, Lemma~\ref{lem:filter} and Markov's inequality give, for every $q\ge1$,
\[
 T_\gamma(r)\le r^{-q}\int_\R|t|^q|W_\gamma(t)|\,dt.
\]
Choose $q=N+2$ and sum $n^{-q}$ over $n\ge\lfloor d\rfloor$.
\end{proof}

\begin{theorem}[conditional continuum LPPL]\label{thm:main}
Assume Hypotheses~\ref{hyp:G} and \ref{hyp:M}.  Then there is a state $\omega_\infty$ on $\A$ such
that, for every bounded interval $O$,
\[
 \lim_{g\in\G}\left\|(\omega_g-\omega_\infty)|_{\A(O)}\right\|=0.
\]
For every $g$ with $g=1$ on $O_d$,
\begin{equation}\label{eq:limit-rate}
 \left\|(\omega_g-\omega_\infty)|_{\A(O)}\right\|
 \le\Psi_\gamma(d),
\end{equation}
and $\Psi_\gamma$ obeys Lemma~\ref{lem:Psitail}.  The restriction
$\omega_\infty|_{\A(O)}$ is normal for every $O$.

If \eqref{eq:physical} holds, $\omega_\infty$ is a ground state of $(\A,\alpha)$ in the spectral
sense of M0.  If \eqref{eq:translation} holds, $\omega_\infty$ is translation invariant.  If $P$
is even and every $\omega_g$ is $\mathbb Z_2$-invariant, then so is $\omega_\infty$.
\end{theorem}

\begin{proof}
Fix $O$ and $\varepsilon>0$.  By Lemma~\ref{lem:Psitail}, choose $d$ with
$\Psi_\gamma(d)<\varepsilon$.  Choose $g_0\in\G$ with $g_0=1$ on $O_d$.  If
$g,g'\succcurlyeq g_0$, then $g=g'=1$ on $O_d$, so Proposition~\ref{prop:path} gives
\[
 \| (\omega_g-\omega_{g'})|_{\A(O)}\|<\varepsilon.
\]
Thus $(\omega_g|_{\A(O)})_{g\in\G}$ is Cauchy in the Banach space $\A(O)^*$.
Its limit is a state $\omega_O$.  The limits are compatible under isotony, so they define a
positive norm-one functional on $\bigcup_O\A(O)$, which extends uniquely by norm continuity to a
state $\omega_\infty$ on $\A$.

To prove \eqref{eq:limit-rate}, first suppose $d\ge d_0$.  Fix $g=1$ on $O_d$ and let $g'$ tend to
infinity through the cofinal tail of cutoffs which are also $1$ on $O_d$.
Proposition~\ref{prop:path} applies to every such $g'$, and norm convergence permits passage to the
limit.  If $0\le d<d_0$, the assertion follows from
$\|\omega_g-\omega_\infty\|\le2\le\Psi_\gamma(d)$.

Each $\omega_g|_{\A(O)}$ is normal.  Normal functionals form a norm-closed subspace of
$\A(O)^*$; explicitly, if $0\le A_i\uparrow A$ in $\A(O)$ and
$\|\omega_g-\omega_\infty\|_{\A(O)^*}<\delta$, then
\[
 0\le\omega_\infty(A)-\omega_\infty(A_i)
 \le2\delta\|A\|+\omega_g(A-A_i).
\]
First let $i$ increase, using normality of $\omega_g$, and then let $\delta\downarrow0$.
Hence $\omega_\infty|_{\A(O)}$ is normal.

Choose an increasing sequence $(\widetilde g_n)_{n\ge1}\subset\G$ such that
$\widetilde g_n=1$ on $(-n,n)$.  For every bounded interval $I$,
\[
 d(\widetilde g_n,I):=
 \dist\!\left(I,\R\setminus\{\widetilde g_n=1\}^{\circ}\right)\longrightarrow\infty,
 \qquad
 \| (\omega_{\widetilde g_n}-\omega_\infty)|_{\A(I)}\|\longrightarrow0.
\]
These are exactly the two hypotheses of M0, Theorem~4.2.  Its exact spectral-passage argument,
using \eqref{eq:physical}, therefore implies that $\omega_\infty$ is an $\alpha$-ground state.

For translation invariance, fix $a\in\R$ and $A\in\A(O)$.  Translation is an order automorphism
of $\G$, so both $g$ and $\tau_ag$ tend cofinally to infinity.  By \eqref{eq:translation},
\[
 \omega_\infty(A)
 =\lim_g\omega_{\tau_ag}(A)
 =\lim_g\omega_g(\tau_{-a}A)
 =\omega_\infty(\tau_{-a}A).
\]
Density extends this from local $A$ to $\A$.  The $\mathbb Z_2$ assertion follows directly by
taking the local norm limit of invariant states.
\end{proof}

\begin{remark}[what the theorem does not prove]\label{rem:limits}
Theorem~\ref{thm:main} proves uniqueness of the limit of the specified cutoff family.  It does not
show that every locally normal ground state of $(\A,\alpha)$ equals $\omega_\infty$.  Identification
with a separately constructed Euclidean vacuum also requires either a direct comparison theorem
or an independent uniqueness result.  The super-polynomial function $\Psi_\gamma$ is an explicit
valid rate furnished by the chosen exact filter; no claim of a sharp $e^{-2md}$ rate is made.
\end{remark}

\section{Remaining obligations}\label{sec:status}

Proposition~\ref{prop:standard} closes exact cutoff propagation and every vector-domain assertion.
Proposition~\ref{prop:FKNpath} closes affine-path differentiability and endpoint continuity under
the already assumed gap.  Theorem~\ref{thm:main} therefore reduces the remaining continuum Cauchy
step to two statements:
\begin{enumerate}
\item prove the full-net uniform physical gap of Hypothesis~\ref{hyp:G} in a precisely specified
      noncritical parameter regime;
\item prove the uniform local second moment in Hypothesis~\ref{hyp:M} by a valid ultraviolet
      estimate.
\end{enumerate}
The discharge audit and explicit reductions are in \texttt{remaining.tex}.  In particular, the
covariance-weighted Poincar\'e induction previously proposed for the moment bound is invalid and is
not used.  A wavelet
resolution requires additional comparison maps, a uniform interacting Lieb--Robinson estimate,
and ultraviolet convergence, and is therefore a separate theorem rather than part of
Theorem~\ref{thm:main}.

\begin{thebibliography}{9}
\bibitem{GJExp}
J.~Glimm and A.~Jaffe, \emph{Quantum Field Theory and Statistical Mechanics: Expositions},
Birkh\"auser, Boston, 1985.

\bibitem{M0}
\emph{Limit points of the finite-volume ground states of $P(\varphi)_2$ are ground states of the
quasi-local dynamics}, Milestone M0, \texttt{../m0/m0.tex}.

\bibitem{M3}
\emph{Clustering, the symmetry-breaking dichotomy, and the reduction of uniqueness},
Milestone M3, \texttt{../m3/m3.tex}.

\bibitem{M4b}
\emph{Discharge audit and corrected proof program for continuum LPPL in $P(\varphi)_2$},
Milestone M4b, \texttt{remaining.tex}.

\bibitem{Simon}
B.~Simon, \emph{The $P(\varphi)_2$ Euclidean (Quantum) Field Theory},
Princeton University Press, Princeton, 1974.
\end{thebibliography}

\end{document}
